\section{The supported local descent}
\label{sec:descent}

This section expands the middle of the argument.  In particular, none of the
implications below is justified by saying that the proof for smooth moduli is
unchanged.  We retain the support needed at the next stage, state every
enlargement in the direction in which it is used, and record its loss.

Fix one nonempty outer localization.  Thus $Q,R,q_0,H,H^*,U,V$ satisfy
\cref{eq:H,eq:UV,eq:Delta1}, with $1\le H^*\ll H$, and put

\begin{equation}\label{eq:g-g0}
 g=(v_1,v_2),\qquad g_0=(q_0,\ell),\qquad
 v_i=gv_i^*,\qquad (v_1^*,v_2^*)=1.
\end{equation}

All coefficient sequences in this section are bounded by
$(\log x)^{O(1)}$.  Such factors, as well as a fixed number of dyadic
partitions, are written as $x^{o(1)}$.  Since $\eta$ is fixed, each of them
is smaller than $x^{\eta/100}$ for sufficiently large $x$.

\subsection{Preparation and the first phase identity}

Choose the witness $d_1\mid r$ before duplicating any variable, write
$r=d_1r_1$, and localize $d_1\asymp\Delta$.  In the notation of
\cref{lem:second-factor},

\begin{equation}\label{eq:Delta-range-descent}
 S_x^{-2}\frac{N}{q_0^2x^{\delta+50\eta}H^2}
 \ll\Delta\ll
 \frac{N}{q_0^2x^{50\eta}H^2},
 \qquad \Delta_1=\Delta x^{-5\eta}.
\end{equation}

For fixed $r_1,u_1,v_1,v_2,q_2,d_0$, let $\Sigma_2$ denote

\begin{align}
\Sigma_2={}&
 \sum_{\substack{d\\(d,m)=1\\d\ {\rm squarefree}}}
 \psi_{\Delta_1}(d-d_0)
 \Bigg|
 \sum_{\substack{h_1,h_2\sim H^*\\h_1v_2\ne h_2v_1}}
 \sum_{n}^{\dagger}
 C(n;d)\psi_N(n)\notag\\[-1mm]
 &\hspace{23mm}\times
 \varphi_1(h_1)\overline{\varphi_2(h_2)}
 \Psi(n,h_1,h_2;d)
 \Bigg|,
 \label{eq:Sigma2}
\end{align}

where

\begin{equation}\label{eq:dagger2}
 m=r_1q_0u_1[v_1,v_2]q_2,\qquad
 (n,dr_1q_0u_1v_1v_2)=1,\qquad
 (n+\ell dr_1,q_0q_2)=1.
\end{equation}

The indicator $C(n;d)$ is supported on at most $g_0$ classes of $n$ modulo
$q_0$ for each class of $d$ modulo $q_0$.

\begin{lemma}[Prepared-factor implication]\label{lem:prepared-factor}
Uniformly in the fixed outer data, the estimate

\begin{equation}\label{eq:Sigma2-target}
 \Sigma_2\ll q_0g_0\Delta Ngx^{-10\eta}
\end{equation}

implies

\begin{equation}\label{eq:Sigma1-target}
 \Sigma_1^\sharp\ll g_0RQNUV^2x^{-4\eta}.
\end{equation}
\end{lemma}

\begin{proof}
Grouping the original $r$-sum according to $r=d_1r_1$, followed by dyadic
localization, gives

\begin{equation}\label{eq:prep-group}
 \Sigma_1^\sharp
 \ll x^{\eta/2}\frac{RUQ}{q_0\Delta}
 \sup_{\Delta,r_1,u_1,q_2}\Upsilon_0.
\end{equation}

This is a counting step: the number of choices of the frozen variables is
$O(x^{\eta/2}RUQ/(q_0\Delta))$ after the usual divisor-weight
enlargement.  No factor of a prescribed size is chosen here; the factor was
already supplied by $d_1(r)$.

Separate the relation $h_1v_2=h_2v_1$.  Since
$H^*\ll H\le x^{-5\eta}V$ by \cref{eq:UV}, the number of solutions to that
relation is $O(x^{o(1)}H^*V)$, and its contribution is

\begin{equation}\label{eq:first-diagonal}
 \Upsilon_{0,=}
 \ll x^{o(1)}\Delta H^*VN
 \ll x^{-5\eta+o(1)}\Delta NV^2.
\end{equation}

For the complementary part, cover $d\asymp\Delta$ by
$O(x^{5\eta})$ shifted smooth intervals of length $\Delta_1$.  After the
outer absolute value has been taken, deleting the original family-membership
indicators and the fixed coefficient of modulus at most one can only increase
the sum.  The remaining squarefree and coprimality conditions are exactly
those in \cref{eq:Sigma2}.  Thus \cref{eq:Sigma2-target} and

\begin{equation}\label{eq:gcd-v-average}
 \sum_{v_1\asymp V}\sum_{v_2\asymp V}(v_1,v_2)
 \ll V^2(\log V)^{O(1)}
\end{equation}

give

\[
 \Upsilon_{0,\ne}
 \ll q_0g_0\Delta NV^2x^{-5\eta+o(1)}.
\]

The right side of \cref{eq:first-diagonal} is smaller than this admissible
majorant.  Insert the result in \cref{eq:prep-group}; the factors $q_0$ and
$\Delta$ cancel, and the powers $x^{\eta/2}$ and
$x^{-5\eta+o(1)}$ leave at least $x^{-4\eta}$.  This proves
\cref{eq:Sigma1-target}.
\end{proof}

\begin{lemma}[CRT phase factorization]\label{lem:phase-factorization}
Under the support conditions in \cref{lem:prime-support}, there are residue
classes $A,B\pmod m$, independent of $n,d,h_1,h_2$, such that $(A,m)=1$
and, with $y=h_1v_2^*-h_2v_1^*$,

\begin{equation}\label{eq:phase-factorization}
 \Psi(n,h_1,h_2;d)
 =e_d\!\left(
   \frac{ay}{nr_1q_0u_1gv_1^*v_2^*q_2}
  \right)
  e_m\!\left(\frac{Ay}{d(n+Bd)}\right).
\end{equation}

One may choose $A,B$ by

\begin{align}
 A&\equiv a\pmod{r_1},& B&\equiv0\pmod{r_1},\notag\\
 A&\equiv b_1\pmod{q_0u_1gv_1^*v_2^*},&
 B&\equiv0\pmod{q_0u_1gv_1^*v_2^*},\notag\\
 A&\equiv b_2\pmod{q_2},& B&\equiv\ell r_1\pmod{q_2}.
 \label{eq:AB-CRT}
\end{align}
\end{lemma}

\begin{proof}
By \cref{lem:prime-support}, the three moduli in
\cref{eq:AB-CRT} are pairwise coprime and their product is $m$.
The Chinese remainder theorem therefore defines $A$ and $B$ uniquely modulo
$m$.  Splitting each additive character of the duplicated phase over the
prime-support decomposition and collecting the $d$-component gives the first
factor in \cref{eq:phase-factorization}; collecting the remaining components
gives the second.  The residue $A$ is a unit because $a,b_1,b_2$ are units on
their respective support factors.  This proves the identity and every
inverse appearing in it is legitimate.
\end{proof}

\subsection{The algebraic \texorpdfstring{$q$}{q}-van der Corput step}

Put

\begin{equation}\label{eq:Yset}
 \mathscr Y=\{h_1v_2^*-h_2v_1^*:h_1,h_2\sim H^*\}.
\end{equation}

For a nonzero $y\in\mathscr Y$, set

\[
 w_1=(y,d),\qquad w_2=(y,m^*),\qquad
 m^*=m^{\lfloor2\log x\rfloor}.
\]

After dyadic localization write $w_1\asymp W_1$ and $y\sim Y$, where

\begin{equation}\label{eq:WY-ranges}
 1\ll W_1\ll\Delta,\qquad
 1\ll |Y|\ll\frac{H^*V}{g}.
\end{equation}

For $w_0\mid w_1$, define

\begin{equation}\label{eq:K3-def}
 \mathscr K_3(n,\widetilde n;y,\widetilde y,d)
 :=e_m\!\left(
 \frac{A\{y(\widetilde n+Bd)-\widetilde y(n+Bd)\}}
 {d(n+Bd)(\widetilde n+Bd)}
 \right).
\end{equation}

\begin{align}
\Sigma_3:={}&
 \sum_{\substack{y,\widetilde y\in\mathscr Y\\
 y,\widetilde y\sim Y\\w_1\mid(y,\widetilde y)\\
 (y,m^*)=(\widetilde y,m^*)=w_2}}
 \Bigg|
 \sum_{\substack{d\\w_0w_1\mid d\\
 (d/w_1,my\widetilde y/w_1^2)=1}}
 \psi_{\Delta_1}(d-d_0)\notag\\[-1mm]
&\hspace{17mm}\times
 \sum_{\substack{n,\widetilde n\\
 y\widetilde n\equiv\widetilde yn\ (d)}}^{\dagger_3}
 C(n;d)C(\widetilde n;d)\psi_N(n)\psi_N(\widetilde n)\notag\\[-1mm]
&\hspace{30mm}\times
 \mathscr K_3(n,\widetilde n;y,\widetilde y,d)
 \Bigg|,
 \label{eq:Sigma3}
\end{align}

where $\dagger_3$ retains

\begin{equation}\label{eq:dagger3}
 (n\widetilde n,w_1r_1q_0u_1v_1v_2)=1,\qquad
 ((n+\ell dr_1)(\widetilde n+\ell dr_1),q_0q_2)=1.
\end{equation}

\begin{lemma}[Supported $q$-van der Corput inequality]
\label{lem:qvdc-supported}
For some admissible $W_1,Y,w_0,w_1,w_2$,

\begin{equation}\label{eq:qvdc-bound}
 \Sigma_2\ll x^{3.5\eta+o(1)}g\Delta\,\Sigma_3^{1/2}.
\end{equation}

Consequently

\begin{equation}\label{eq:Sigma3-target}
 \sup\Sigma_3\ll q_0^2g_0^2N^2x^{-27\eta}
\end{equation}

implies \cref{eq:Sigma2-target}.
\end{lemma}

\begin{proof}
On a fixed $d$, partition the inner sum by
$yn^{-1}\pmod d$.  In \cref{eq:phase-factorization} the $e_d$ factor then
depends only on that residue class and disappears after absolute values.
Retain $d$ squarefree until $w_1=(y,d)$ is fixed.  It follows that $w_1$ is
squarefree and $(w_1,m)=1$.

Only $x^{o(1)}$ values of $w_2$ occur.  Indeed, a nonempty block has
$w_2\le |y|\le x^{C_1}$ for a fixed $C_1$.  If
$m=P_1\cdots P_r$ with $P_1<\cdots<P_r$, replace each $P_j$ by the
$j$-th prime.  This injects the possible divisors of $m^*$ below $x^{C_1}$
into the $p_r$-smooth integers below that bound.  Since $m\le x^{C_*}$ is
squarefree, $r=O(\log x/\log\log x)$ and $p_r=O(\log x)$.  Rankin's
bound with $\sigma=1/\log p_r$ gives

\[
 \#\{w_2:w_2\mid m^*,\ w_2\le x^{C_1}\}
 \le x^{C_1\sigma}\prod_{p\le p_r}(1-p^{-\sigma})^{-1}
 =x^{o(1)}.
\]

Dyadic decomposition in $w_1$ and $y$ and the sum over residue classes cost
at most $x^{2\eta+o(1)}W_1$.  Enlarge the squarefree $d$-sum to the
conditions $w_1\mid d$ and $(d/w_1,w_1)=1$.  The congruence
$yn^{-1}\pmod d$ can then be inverted modulo $d/w_1$.  Replacing
$(n,d)=1$ by the weaker $(n,w_1)=1$ is a nonnegative enlargement.

Cauchy--Schwarz in the residue class and $d$ variables contributes
\[
 x^{\eta/2}\Delta/W_1.
\]
The preceding factor $W_1$ cancels.  The accumulated factor is
$x^{2.5\eta+o(1)}\Delta$.  The two copied
congruences combine to

\[
 y\widetilde n\equiv\widetilde yn\pmod d.
\]

For fixed $y$, the equation $y=h_1v_2^*-h_2v_1^*$ has at most $O(g)$
solutions in $h_1,h_2\sim H^*$: all solutions differ by
$(kv_1^*,kv_2^*)$.  Applying this to $y$ and $\widetilde y$ and then taking
the square root contributes $gx^{\eta/2+o(1)}$.  Finally

\[
 1_{(d/w_1,w_1)=1}
 =\sum_{w_0\mid(d/w_1,w_1)}\mu(w_0)
\]

and the divisor count of the squarefree $w_1$ contributes another
$x^{\eta/2+o(1)}$.  The resulting inner expression is exactly
\cref{eq:Sigma3}, proving \cref{eq:qvdc-bound}.  Taking the square root of
\cref{eq:Sigma3-target} leaves
$q_0g_0Nx^{-13.5\eta}$; multiplication by the prefactor in
\cref{eq:qvdc-bound} gives \cref{eq:Sigma2-target}.
\end{proof}

\subsection{Removal of diagonals and large gcd blocks}

Let

\begin{equation}\label{eq:tT}
 t=(w_2,m),\qquad
 T=\max\{t^{-1},H^{-1}\}
 \frac{x^{\delta+100\eta}H^2N}{g\Delta_1}.
\end{equation}

For fixed $y,\widetilde y$, let $\Sigma_4(y,\widetilde y)$ be the inner
absolute value in \cref{eq:Sigma3}, restricted further by

\begin{equation}\label{eq:retained-gcd}
 \left(
 \frac{y(\widetilde n+Bd)-\widetilde y(n+Bd)}d,m
 \right)\le T.
\end{equation}

\begin{lemma}[Diagonal and large-gcd removal]\label{lem:diagonal-removal}
One has

\begin{align}
 \Sigma_3\ll{}&
 \max\left\{
 \frac{q_0^2x^{\delta+10\eta}H^3}{t},
 \frac{H^4}{t}
 \right\}
 \sup_{y,\widetilde y}\Sigma_4(y,\widetilde y)
 +x^{-47\eta}N^2.
 \label{eq:diagonal-removal}
\end{align}
\end{lemma}

\begin{proof}
By \cref{eq:WY-ranges,eq:UV},

\begin{equation}\label{eq:Y-corrected}
 |Y|\ll\frac{HV}{g}
 \ll\frac{q_0x^{\delta+5\eta}H^2}{g}.
\end{equation}

If a divisor $q$ of the numerator in \cref{eq:retained-gcd} exceeds $T$,
then \cref{eq:Delta1,eq:tT,eq:Y-corrected} give

\begin{equation}\label{eq:Y-over-q}
 \frac{|Y|}{q}
 \ll q_0^{-1}x^{-150\eta}
 \frac{\min\{t,H\}}{H^2}.
\end{equation}

We count the discarded tuples in four disjoint cases.  If
$y=\widetilde y$ and $n=\widetilde n$, their number is

\begin{equation}\label{eq:diag-1}
 O(\Delta H^2N)
 \ll q_0^{-2}x^{-50\eta}N^2.
\end{equation}

If $y=\widetilde y$ but $n\ne\widetilde n$, the congruence with a fixed
$q>T$ and \cref{eq:Y-over-q} give

\begin{equation}\label{eq:diag-2}
 \max\{H^2/t,H\}N^2\frac{|Y|}{q}
 \ll q_0^{-1}x^{-150\eta}N^2.
\end{equation}

If $n=\widetilde n$ and $y\ne\widetilde y$, split according to
$(y-\widetilde y,dq)$.  Divisor expansion bounds the resulting two terms by

\begin{equation}\label{eq:diag-3}
 x^{\delta-144\eta+o(1)}N,
 \qquad
 x^{2\delta-39\eta+o(1)}H^2N.
\end{equation}

The inequalities following from \cref{eq:S52-a,eq:S52-c}, namely
$\gamma>\delta$ and $\gamma>8\omega+4\delta$, together with
$H\ll x^{4\omega+\delta+7\eta}/q_0$, make both terms
$O(x^{-48\eta}N^2)$ after reducing $\eta_0$.

In the fully off-diagonal case, fix $d,\widetilde n,y-\widetilde y$.  The
defining equation leaves $O(q_0\min\{t,H\})$ possible quotient values and
only $x^{o(1)}$ factorizations for each.  Hence this contribution is

\begin{equation}\label{eq:diag-4}
 x^{o(1)}q_0\max\{H^2/t,H\}\Delta N\min\{t,H\}
 \ll q_0^{-1}x^{-50\eta+o(1)}N^2.
\end{equation}

The number of possible gcd divisors is $x^{o(1)}$, so
\cref{eq:diag-1,eq:diag-2,eq:diag-3,eq:diag-4} sum to the error in
\cref{eq:diagonal-removal}.

It remains to count retained phase pairs.  The representation through
$h_1,h_2$ gives

\[
 \#(y,\widetilde y)\ll\max\{H^4/t^2,H^2\},
\]

whereas the interval bound gives
$\#(y,\widetilde y)\ll Y^2/t^2$.  If $t\le x^\delta H$, use the first
estimate; if $t>x^\delta H$, use the second and
\cref{eq:Y-corrected}.  In both cases,

\begin{equation}\label{eq:retained-pairs}
 \#(y,\widetilde y)
 \ll\max\left\{
 \frac{q_0^2x^{\delta+10\eta}H^3}{t},
 \frac{H^4}{t}
 \right\}.
\end{equation}

Multiplying the supremum of $\Sigma_4$ by this count completes the proof.
\end{proof}

\subsection{Relabelling the retained blocks}

Write

\begin{equation}\label{eq:relabel}
 z_1=w_0w_1,\qquad y=w_1\lambda,\qquad
 \widetilde y=w_1\widetilde\lambda,\qquad d=z_1e,
 \qquad \Lambda=Y/w_1.
\end{equation}

An empty or incompatible block is discarded.  In every remaining block,

\begin{gather}
 z_1\ll x^{5\eta}\Delta_1,\qquad (z_1,m)=1,\qquad
 w_1\mid z_1,\qquad
 (\lambda,m^*)=(\widetilde\lambda,m^*)=w_2,
 \label{eq:Z5-arithmetic}\\
 H\ll q_0^{-1}x^{4\omega+\delta+7\eta},\qquad
 x^{-4\eta-\delta}N\ll R\ll x^{-2\eta}N,
 \label{eq:Z5-ranges-a}\\
 x^{1/2-\eta}\ll QR\ll x^{1/2+2\omega+\eta},\qquad
 x^{5\eta}H\ll v_i\ll q_0x^{\delta+5\eta}H,
 \label{eq:Z5-ranges-b}\\
 \frac{RQ^2H}{q_0g\Delta_1}\ll m
 \ll\frac{x^\delta RQ^2H}{g\Delta_1},\qquad
 1\ll|\Lambda|\ll
 \frac{q_0x^{\delta+5\eta}H^2}{w_1g}.
 \label{eq:Z5-ranges-c}
\end{gather}

The apparent asymmetry in the two bounds for $m$ is genuine: the lower
endpoint contains the common factor removed from the first prescribed
divisor, whereas the upper endpoint uses only its ambient dyadic range.

Replace $e$ again by $d$, and set

\begin{equation}\label{eq:relabel-rest}
 \begin{aligned}
 l&=\ell z_1r_1,&
 c_1&=r_1q_0u_1[v_1,v_2],& c_2&=q_0q_2,\\
 A&\leftarrow Aw_1/z_1,&
 B&\leftarrow Bz_1.
 \end{aligned}
\end{equation}

All quotients in \cref{eq:relabel-rest} are units modulo $m$.
The retained sum becomes

\begin{align}
\Sigma_5={}&\Bigg|
 \sum_{(d,m\lambda\widetilde\lambda)=1}
 \psi_{\Delta_1}(z_1d-d_0)
 \sum_{n,\widetilde n}^{\dagger_5}
 C(n;d)C(\widetilde n;d)\psi_N(n)\psi_N(\widetilde n)
 \notag\\[-1mm]
&\hspace{12mm}\times
 e_m\!\left(
 \frac{A\{\lambda(\widetilde n+Bd)-
 \widetilde\lambda(n+Bd)\}}
 {d(n+Bd)(\widetilde n+Bd)}
 \right)\Bigg|,
 \label{eq:Sigma5}
\end{align}

where $\dagger_5$ means

\begin{gather}
 (n\widetilde n,w_1c_1)=1,\qquad
 ((n+ld)(\widetilde n+ld),c_2)=1,
 \label{eq:dagger5-a}\\
 \lambda\widetilde n\equiv\widetilde\lambda n
 \pmod{z_1d/w_1},\qquad
 \left(
 \frac{\lambda(\widetilde n+Bd)-\widetilde\lambda(n+Bd)}d,m
 \right)\le T.
 \label{eq:dagger5-b}
\end{gather}

Equations \cref{eq:relabel,eq:relabel-rest} are reversible on each nonempty
block, so $\Sigma_4=\Sigma_5$.  Combining
\cref{eq:diagonal-removal,eq:Sigma3-target} shows that it is enough to prove

\begin{equation}\label{eq:Sigma5-target}
 \Sigma_5\ll
 \min\left\{\frac{t}{q_0^2x^{\delta+10\eta}H^3},
              \frac{t}{H^4}\right\}
 q_0^2g_0^2N^2x^{-27\eta}.
\end{equation}

\subsection{First elimination: the copied variable}

The congruence in \cref{eq:dagger5-b} is equivalent to

\begin{equation}\label{eq:ntilde-param}
 \widetilde n=
 \frac{\widetilde\lambda n+k(z_1/w_1)d}{\lambda}
\end{equation}

for an integer $k$.  The supports of the two $N$-scale profiles and of the
$\Delta_1$-scale profile imply

\begin{equation}\label{eq:k-range}
 |k|\ll K_0:=\frac{w_1|\Lambda|N}{x^{5\eta}\Delta_1}.
\end{equation}

Put

\begin{equation}\label{eq:S-k}
 S_k=\frac{z_1}{w_1}k+(\lambda-\widetilde\lambda)B.
\end{equation}

For shifted smooth profiles $\psi_N^*,\psi_{\Delta_1}^*$ and residue
classes $d_*,n_*\pmod{q_0}$, define

\begin{align}
\Sigma_6(k):={}&\Bigg|
 \sum_{\substack{d\equiv d_*\ (q_0)\\
 (d,m\lambda\widetilde\lambda)=1}}
 \psi_{\Delta_1}^*(z_1d-d_0)
 \sum_{\substack{n\equiv n_*\ (q_0)\\
 \lambda\mid\widetilde\lambda n+k(z_1/w_1)d}}^{\dagger_6}
 \psi_N^*(n)
 \notag\\[-1mm]
&\hspace{12mm}\times
 e_m\!\left(
 \frac{AS_k}
 {(n+Bd)
  \{\lambda^{-1}(\widetilde\lambda n+k(z_1/w_1)d)+Bd\}}
 \right)\Bigg|,
 \label{eq:Sigma6}
\end{align}

where $\dagger_6$ retains the four coprimality conditions obtained by
substituting \cref{eq:ntilde-param} into \cref{eq:dagger5-a}.

\begin{lemma}[Elimination of $\widetilde n$]\label{lem:remove-ntilde}
If nonnegative weights $\xi_k$ satisfy

\begin{equation}\label{eq:xi-mass}
 \sum_{\substack{|k|\ll K_0\\w_2\mid k\\(S_k,m)\le T}}\xi_k\ll1
\end{equation}

and, uniformly in the profiles and residue classes,

\begin{equation}\label{eq:Sigma6-target}
 \Sigma_6(k)\ll \xi_k
 \min\left\{\frac{t}{q_0^2x^{\delta+10\eta}H^3},
              \frac{t}{H^4}\right\}
 q_0g_0N^2x^{-27\eta},
\end{equation}

then \cref{eq:Sigma5-target} holds.
\end{lemma}

\begin{proof}
Substitution of \cref{eq:ntilde-param} in the phase gives exactly
\cref{eq:S-k,eq:Sigma6}.  For the copied smooth weight, write

\begin{align}
 &\psi_N\!\left(
 \frac{\widetilde\lambda n+k(z_1/w_1)d}{\lambda}
 \right)\notag\\
 &\quad=
 \sum_{j=0}^{J}\frac1{j!}
 \left(\frac{(k/w_1)(z_1d-d_0)}{\lambda N}\right)^j
 \psi^{(j)}\!\left(
 \frac{\widetilde\lambda n+(k/w_1)d_0}{\lambda N}
 \right)+O(x^{-100}),
 \label{eq:Taylor-one}
\end{align}

where $J=\lceil20/\eta\rceil$.  On the support of the sum, the expansion
parameter is at most $x^{-5\eta}$ by \cref{eq:k-range}.  Each summand in
\cref{eq:Taylor-one} is a product of one smooth $N$-profile in $n$ and one
shifted smooth $\Delta_1$-profile in $z_1d-d_0$.  The remainder contributes
$O(x^{-90})$ after trivial summation.

For fixed $k$, the two indicators $C(n;d)$ and
$C(\widetilde n;d)$ are constant after $d,n$ are fixed modulo $q_0$.
There are at most $q_0g_0$ contributing pairs of residue classes.  Moreover,
the divisibility condition in \cref{eq:Sigma6} forces $w_2\mid k$: both
$\lambda$ and $\widetilde\lambda$ are divisible by $w_2$, whereas
$(z_1d/w_1,w_2)=1$.  Summing \cref{eq:Sigma6-target} over $k$ and using
\cref{eq:xi-mass} supplies the missing factor $q_0g_0$, and proves
\cref{eq:Sigma5-target}.  The Taylor remainder is smaller than the target by
the fixed parameter reserve.
\end{proof}

\subsection{Second elimination: the divisibility by
\texorpdfstring{$\lambda$}{lambda}}

Let

\begin{equation}\label{eq:lambda-star}
 \lambda^*=(|\lambda|,|\widetilde\lambda|),\qquad
 a_\lambda=\lambda/\lambda^*,\qquad
 b_\lambda=\widetilde\lambda/\lambda^*.
\end{equation}

By \cref{lem:valuation}, $a_\lambda$ and $b_\lambda$ are units modulo $m$.
We write $\bar a_\lambda,\bar b_\lambda$ for their inverses modulo $m$.
Likewise, $\overline{\widetilde\lambda/w_2}$ denotes the inverse of
$\widetilde\lambda/w_2$ modulo $m$.

If $\lambda^*\nmid(z_1/w_1)k$, then \cref{eq:Sigma6} is empty.  Otherwise
there are residue classes $F_k,G_k$ such that

\begin{equation}\label{eq:n-n1}
 n=a_\lambda n_1+F_kd,\qquad
 \widetilde n=b_\lambda n_1+G_kd.
\end{equation}

They satisfy

\begin{equation}\label{eq:FG-relation}
 \bar b_\lambda(G_k+B)-\bar a_\lambda(F_k+B)
 \equiv
 \bar a_\lambda\,\overline{\widetilde\lambda/w_2}\,
 \frac{S_k}{w_2}\pmod m.
\end{equation}

Define the genuine short scale

\begin{equation}\label{eq:Delta-star}
 \Delta^*=\min\left\{\frac{N}{|\Lambda|x^{5\eta}},\Delta_1\right\}.
\end{equation}

For shifted profiles at scales $\Delta^*/z_1$ and
$|\bar a_\lambda|_{\rm sc}N$ (the latter notation means the real scale
$|\lambda^*/\lambda|N$), let

\begin{align}
\Sigma_7(k):={}&\Bigg|
 \sum_{\substack{d\equiv d_*\ (q_0)\\
 (d,m\lambda\widetilde\lambda)=1}}
 \psi_{\Delta^*/z_1}(d)
 \sum_{n_1}^{\dagger_7}\psi_{|\lambda^*/\lambda|N}(n_1)
 \notag\\[-1mm]
&\quad\times e_m\!\left(
 \frac{\bar a_\lambda\bar b_\lambda AS_k}
 {(n_1+\bar a_\lambda(F_k+B)d)
  (n_1+\bar b_\lambda(G_k+B)d)}
 \right)\Bigg|.
 \label{eq:Sigma7}
\end{align}

Here $\dagger_7$ consists of the congruence
$a_\lambda n_1+F_kd\equiv n_*\pmod{q_0}$ and the four coprimality
conditions obtained from \cref{eq:n-n1}.

\begin{lemma}[Divisibility elimination]\label{lem:divisibility-elimination}
For every retained $k$,

\begin{equation}\label{eq:Sigma6-to-7}
 \Sigma_6(k)
 \ll\frac{\Delta_1}{\Delta^*}
 \sup\bigl(\Sigma_7(k)+O(x^{-90})\bigr),
\end{equation}

where the supremum is over the shifted profiles and the compatible choices
of $F_k,G_k$.
\end{lemma}

\begin{proof}
Dividing the condition
$\lambda\mid\widetilde\lambda n+(z_1/w_1)kd$ by $\lambda^*$ gives one
linear class for $n$ modulo $a_\lambda$ because
$(a_\lambda,b_\lambda)=1$.  Choosing its representative $F_k$ and
substituting gives \cref{eq:n-n1}; the equality defining $G_k$ then gives
\cref{eq:FG-relation} by direct subtraction.

Use a smooth partition of unity to split the original $d$-support into
$O(\Delta_1/\Delta^*)$ intervals of length $\Delta^*/z_1$.  On one such
interval, centered at $x_0$, the only mixed profile is
$\psi_N^*(a_\lambda n_1+F_kd)$.  Its normalized displacement satisfies

\[
 \frac{|F_k|\,|d-x_0|}{N}
 \ll \frac{|\lambda|}{\lambda^*}
       \frac{\Delta^*}{z_1N}
 \ll x^{-5\eta}.
\]

Taylor expansion to order $J$ separates this profile into products of a
shifted $n_1$-profile and a shifted $d$-profile; the remainder is
$O(x^{-90})$ after summation.  The number of partition blocks is the displayed
factor in \cref{eq:Sigma6-to-7}.  No replacement of $\Delta^*$ by
$\Delta_1$ has been made.
\end{proof}

\subsection{Third elimination: all coprimality conditions}

For $q_3\mid m/q_0$, residue classes $d^*,n^*\pmod{q_0q_3}$, shifted
profiles with $N_2\le N$ and $\Delta_2\le\Delta^*$, and units
$A_1,A_2\pmod m$, define

\begin{align}
\Sigma_8(k):={}&\Bigg|
 \sum_{\substack{d\equiv d^*\ (q_0q_3)}}
 \sum_{\substack{n\equiv n^*\ (q_0q_3)}}
 \psi_{\Delta_2/z_1}(d)\psi_{N_2}(n)\notag\\[-1mm]
&\hspace{27mm}\times
 e_m\!\left(
 \frac{w_2A_1L_1}
 {(n+B_1d)(n+(B_1+L_1)d)}
 \right)\Bigg|,
 \label{eq:Sigma8}
\end{align}

where

\begin{equation}\label{eq:L1}
 L_1\equiv A_2\frac{S_k}{w_2}\pmod m,
 \qquad (A_1A_2,m)=1.
\end{equation}

\begin{lemma}[Coprimality elimination]\label{lem:coprimality-elimination}
For every compatible choice in \cref{eq:Sigma7},

\begin{equation}\label{eq:Sigma7-to-8}
 \Sigma_7(k)
 \ll x^{2\eta+o(1)}q_3\sup\Sigma_8(k)+O(x^{-89}).
\end{equation}
\end{lemma}

\begin{proof}
We give the full sequence of finite transformations.  First apply Mobius
inversion to

\[
 (d,(\lambda/w_2)(\widetilde\lambda/w_2))=1,
 \quad(a_\lambda n_1+F_kd,w_1)=1,
 \quad(b_\lambda n_1+G_kd,w_1)=1.
\]

This introduces

\begin{equation}\label{eq:f123}
 f_1\mid(\lambda/w_2)(\widetilde\lambda/w_2),\qquad
 f_2\mid w_1,\qquad f_3\mid w_1,
\end{equation}

and costs $x^{\eta+o(1)}$.  Put

\[
 f_2^*=(f_2,a_\lambda),\qquad
 f_3^*=(f_3,b_\lambda).
\]

Solvability of the two resulting linear congruences first forces a divisor
$f_{23}$ of $d/f_1$.  On common prime factors at which their required
residues disagree, it forces a further squarefree divisor $f_4$.  The
Chinese remainder theorem then gives the exact parametrization

\begin{equation}\label{eq:f-param}
 \begin{aligned}
 d&=f^*d_3,&
 n_1&=w_3n_3+H_kf_4d_3,\\
 f^*&=f_1f_{23}f_4,&
 w_3&=\lcm(f_2/f_2^*,f_3/f_3^*).
 \end{aligned}
\end{equation}

The supports in \cref{eq:f123} and \cref{lem:valuation} give
$(f^*w_3,m)=1$ and $w_3\mid w_1$.  Hence every division in
\cref{eq:f-param} is valid modulo $m$.

The profile in $w_3n_3+H_kf_4d_3$ is the last mixed profile.  On the
$d_3$-support,

\[
 \frac{|H_kf_4|\,\Delta^*}{z_1f^*|\lambda^*/\lambda|N}
 \ll \frac{w_1|\Lambda|\Delta^*}{z_1N}
 \le x^{-5\eta}.
\]

Taylor expansion through order $J$ therefore separates it, with total
remainder $O(x^{-89})$.

Apply Mobius inversion a second time to the five remaining conditions.  It
introduces

\begin{equation}\label{eq:e01234}
 e_0\mid m,\qquad e_1,e_2\mid c_1,\qquad e_3,e_4\mid c_2
\end{equation}

and a further $x^{\eta+o(1)}$ loss.  Let

\[
 g_1=\lcm(q_0,e_0,e_1,e_2,e_3,e_4)=q_0q_3.
\]

Because $q_0,c_1,c_2\mid m$ and $m$ is squarefree,
$q_3\mid m/q_0$.  The congruences for $d_3,n_3$ are either incompatible,
in which case the block is empty, or determine at most

\[
 \frac{g_1}{\lcm(q_0,e_0)}
 \frac{g_1}{\lcm(q_0,e_1,e_2,e_3,e_4)}
 \le\frac{g_1}{q_0}=q_3
\]

pairs of classes modulo $g_1$.  On each pair, collect the two denominator
slopes into $B_1$ and $B_1+L_1$.  The relation
\cref{eq:FG-relation} gives precisely \cref{eq:L1}; multiplying by units
produces $A_1,A_2$.  The resulting sum is \cref{eq:Sigma8}.  Multiplying
the two Mobius costs and the number of residue pairs proves
\cref{eq:Sigma7-to-8}.
\end{proof}

\subsection{Uniformity of the descendant profiles}

\begin{lemma}[Finite-jet closure]\label{lem:finite-jet}
Let $J=\lceil20/\eta\rceil$ and $J_*=3J+2$.  Suppose every input profile
has fixed compact support and satisfies

\begin{equation}\label{eq:input-jets}
 \|\psi^{(r)}\|_1+\|\psi^{(r)}\|_\infty
 \le B_r(\log x)^{A_r},\qquad 0\le r\le J_*.
\end{equation}

Then every terminal profile in \cref{eq:Sigma8} belongs to one common
shifted $C^2$ profile class, independent of all descendant choices.
Moreover

\begin{equation}\label{eq:terminal-polynomial-size}
 m\le x^7,\qquad N_2\le x^7,\qquad \Delta_2/z_1\le x^7
\end{equation}

for sufficiently large $x$.
\end{lemma}

\begin{proof}
Each of the three Taylor expansions differentiates an input profile at most
$J$ times.  Two terminal derivatives therefore require at most
$3J+2=J_*$ input derivatives.  Products, fixed partitions, and affine
changes of variable preserve the compact support and yield uniform $C^2$
seminorms by Leibniz' rule and the chain rule.  Every normalized Taylor
increment is at most $x^{-5\eta}$, so

\[
 x^{-5\eta(J+1)}(\log x)^{O_\eta(1)}=O(x^{-100}).
\]

There are at most $(J+1)^3$ descendants, a number independent of $x$.

For the size assertion, each of
$r_1,q_0,u_1,v_1,v_2,q_2$ is $O(x)$ and

\[
 m=r_1q_0u_1[v_1,v_2]q_2\le x^6.
\]

Also $N_2\le N\le x$ and
$\Delta_2/z_1\le\Delta^*\le\Delta_1\le\Delta\le x$.
The relaxed common exponent seven leaves room for endpoint constants.
\end{proof}

The lemmas in this section form a one-way chain

\begin{equation}\label{eq:descent-chain}
 \Sigma_8\Longrightarrow\Sigma_7\Longrightarrow\Sigma_6
 \Longrightarrow\Sigma_5\Longrightarrow\Sigma_4
 \Longrightarrow\Sigma_3\Longrightarrow\Sigma_2
 \Longrightarrow\Sigma_1^\sharp.
\end{equation}

Every arrow denotes an upper-bound implication, not an asserted equality.
Restrictions are deleted only after an absolute value makes enlargement
legitimate, and no deleted restriction is later restored.
